\documentclass[11pt]{article}
\usepackage{amsmath,amssymb,amsthm,hyperref}
\newtheorem{theorem}{Theorem}
\newtheorem{lemma}[theorem]{Lemma}
\newtheorem{definition}[theorem]{Definition}
\newcommand{\V}{\mathcal V}
\newcommand{\W}{\mathcal W}
\newcommand{\Aug}{\operatorname{AugFam}}
\newcommand{\id}{\mathrm{id}}
\title{Directed distributivity and a conservative native coefficient envelope}
\date{16 November 2026}
\begin{document}\maketitle
\paragraph{Verification state.} The exact proof text has passed two independent written reviews. This document weakens the invertibility requirement on an independent coefficient input, enlarging the coefficient class for the original occurrence interaction. It does not construct general feedback, joint-scope limits, a free feedback completion, or infinite executions. The construction uses established convolution machinery; no novelty or proof-assistant certification is claimed.

\section{Question and original presentation}
The original independent interaction criterion uses a symmetric rig with invertible distributors. Native maps between finite augmented families aggregate whole coefficient fibres; additive sum is not assumed a categorical coproduct. This matters when extending the source: replacing its coefficients by a convenient probability model can change its actual maps and augmentation semantics. The question here is whether a larger coefficient category can retain the original source fully faithfully while still supporting its native dependent and occurrence calculus.

The answer below uses directed distributors in the contraction orientation. The context and content presentations remain the original ones. The content input is unchanged: an already valid contentad over sets with the specified relative cartesian polynomial occurrence display, cartesian unit/accumulation, and Beck--Chevalley comparisons of the original criterion. No new mixed crossing is assumed.

\section{Independent directed coefficient data}
\begin{definition}[Coherent finite contractions]
A contraction coefficient category has two symmetric monoidal structures $(\oplus,0)$ and $(\otimes,I)$ and natural finite-indexed comparisons
\[
 D:\bigoplus_{(x_i)\in\prod_i X_i}\bigotimes_i V_{i,x_i}
 \longrightarrow\bigotimes_i\bigoplus_{x_i\in X_i}V_{i,x_i}.
\]
These include empty index sets. Unary $D$, the empty-tensor $D$, and $D$ with every $X_i$ singleton are identities up to the chosen structural unit comparisons. The data include the nullary comparisons $0\to0\otimes W$ and $0\to W\otimes0$. Require coherence with finite partitions and their iterated substitution, additive rebracketing and permutations, successive contraction in different tensor variables, and the tensor associator, unitors and symmetry. All comparisons are natural in arbitrary coefficient arrows. Thus contracting a fixed tagged sum-of-tensor expression by either order of partitioning and grouping gives the same map whenever the source, target, and tagged occurrence assignment agree. Unit and zero cases are included in these equations.

Equivalently for the examples used below, take the opposite of a coherent symmetric rig with possibly noninvertible expansion distributors in the sense of the directed distributivity data of Garner--L\'opez Franco, with its multiplicative associativity, unit and symmetry compatibilities. The finite-indexed formulation makes explicit the coherence used by the proof; it asserts no weakest-axiom result. Contractions are never assumed invertible.
\end{definition}
The axioms concern only the coefficient category, sums, tensors, and their structural maps. They mention neither context grades, content grades, nor a mixed interaction law. Original coherent symmetric rigs satisfy them by inverse distributivity. The co-presheaf construction below supplies potentially noninvertible examples.

\section{The full native context calculus}
An object of $\Aug(\V)$ is a finite carrier $X$, coefficients $V_x$, and one augmentation $\epsilon:\bigoplus_x V_x\to I$. A map is a carrier function $f:X\to X'$ and actual maps
\[
 \alpha_{x'}:\bigoplus_{fx=x'}V_x\longrightarrow V'_{x'}
\]
whose partitioned sum $A_\alpha$ preserves the augmentation. Composition partitions the source by successive fibres and composes the given aggregation maps. No directed distributor is involved in defining this category. The unit $J=(1,I,\id)$ is terminal.

For $P=(X,V,\epsilon)$ and conditional $Q_x=(Y_x,W_x,\delta_x)$, define $P\star Q$ with carrier $(x,y)$ and coefficient $V_x\otimes W_{xy}$. Its augmentation is
\[
 \bigoplus_{x,y}V_x\otimes W_{xy}
 \xrightarrow{\oplus_xD}\bigoplus_x\left(V_x\otimes\bigoplus_yW_{xy}\right)
 \xrightarrow{\oplus_x(\id\otimes\delta_x)}\bigoplus_x(V_x\otimes I)
 \cong\bigoplus_xV_x\xrightarrow{\epsilon}I.
\]
For an outer map $(f,\alpha)$ and conditional maps $(g_x,\beta_x):Q_x\to Q'_{fx}$, the coefficient map at $(x',y')$ is
\begin{align*}
 &\bigoplus_{fx=x'}\bigoplus_{g_xy=y'}V_x\otimes W_{xy}
 \xrightarrow{\oplus_xD}
 \bigoplus_{fx=x'}\left(V_x\otimes\bigoplus_{g_xy=y'}W_{xy}\right)\\
 &\xrightarrow{\oplus_x(\id\otimes\beta_{x,y'})}
 \bigoplus_{fx=x'}(V_x\otimes W'_{x'y'})
 \xrightarrow{D}\left(\bigoplus_{fx=x'}V_x\right)\otimes W'_{x'y'}
 \xrightarrow{\alpha_{x'}\otimes\id}V'_{x'}\otimes W'_{x'y'}.
\end{align*}
Thus conditional maps are applied before an outer fibre is merged. Every displayed distributor is a contraction in the declared direction.

\begin{lemma}[Dependent functoriality]
These maps give the original dependent context accumulation functor, preserving augmentations, with invertible associativity and unit comparisons.
\end{lemma}
\begin{proof}
Augmentation preservation is checked from the fully indexed source sum: use contraction naturality and partition coherence to place each continuation augmentation immediately after its corresponding continuation map. Their composite is the source continuation augmentation. The remaining outer composite is $\epsilon' A_\alpha=\epsilon$. This uses no reverse contraction.

For two successive dependent arrows, fix the final target fibre and keep all intermediate indices in the source sum. Partition-substitution coherence compares direct contraction with contraction in stages. The two composites apply the same first conditional maps, second conditional maps, first outer aggregation, and second outer aggregation. In the route with the first outer aggregation earlier, the subsequent conditional map is constant on that merged fibre; tensor bifunctoriality and contraction naturality move it past. No map depending on an unretained source index is moved past a merge.

For an empty source fibre, retain its nullary prefix $\eta_W:0\to0\otimes W$. Naturality says $(0\otimes b)\eta_W=\eta_{W'}$. Consequently, for any actual $\alpha:0\to V'$, one has $(\alpha\otimes b)\eta_W=(\alpha\otimes\id)\eta_{W'}$. This is the required interchange after the prefix; it does not assert $0\otimes b=\id_0$ or cancel $\eta_W$. Multiple empty intermediate fibres use the declared partition/nullary coherence. Nonunique creation maps from $0$ remain present.

Identities use singleton contractions. The associator has carrier $((x,y),z)\mapsto(x,(y,z))$ and coefficient the tensor associator. Both augmentation composites, from the same indexed sum, contract inward and apply the innermost, middle, then outer augmentation. The contraction/tensor compatibility equations equate them. The native-map comparison uses the same continuation-before-parent argument. Pentagon and triangle are the tensor structural equations with the same addressed carrier bijections.
\end{proof}

For a finite addressed family $(P_i)$, its occurrence tensor has carrier $\prod_i X_i$, coefficient $\bigotimes_iV_{i,x_i}$, and augmentation
\[
 \bigoplus_{(x_i)}\bigotimes_iV_{i,x_i}
 \xrightarrow D\bigotimes_i\bigoplus_{x_i}V_{i,x_i}
 \xrightarrow{\otimes_i\epsilon_i}\bigotimes_iI\cong I.
\]
Native maps contract the sum over each joint inverse image and then apply $\bigotimes_i\alpha_{i,x'_i}$.

\begin{lemma}[Occurrence functoriality and native Fubini]
The occurrence tensor is functorial on all aggregated native maps. There is a coherent native isomorphism
\[
 \boxtimes_i(P_i\star Q_i)\cong
 (\boxtimes_iP_i)\star\bigl((x_i)\mapsto\boxtimes_iQ_{i,x_i}\bigr),
\]
whose coefficient is the addressed tensor shuffle. It is natural in arbitrary native outer and conditional maps.
\end{lemma}
\begin{proof}
Tensor-map functoriality compares the contraction of a twice-partitioned joint fibre with successive contractions; this is finite-partition compatibility followed by tensor bifunctoriality. Augmentation preservation follows by naturality and each input augmentation equation.

For Fubini augmentation, start at the sum over all $(x_i,y_i)$. On either side, contraction coherence gives the same contraction toward the grouped outer factors and grouped conditional sums, followed by continuation augmentations and then outer augmentations. Tensor symmetry merely permutes the tagged factors. For map naturality, keep every fibre index visible: contract, apply each $\beta_{i,x_i}$, contract its outer fibre, and apply $\alpha_{i,x'_i}$. Both paths give this composite up to the same tensor shuffle. Empty cases retain the nullary prefixes as above.

The coefficient shuffle and its inverse exist because $\otimes$ is symmetric monoidal, not because $D$ is invertible. Having verified augmentation preservation for the forward shuffle, the inverse also preserves augmentation by composing that equality with the inverse. Compatibility with dependent association and occurrence regrouping compares the same tensor structural maps and the same forward contraction diagrams. No equality is deduced by cancelling a contraction.
\end{proof}

\begin{theorem}[Directed coefficient criterion]
With coherent finite contractions as above and the unchanged original independent content hypotheses, the original category-valued context/content modules admit the original retained-cell pseudodistributive crossing. All native aggregation maps, content maps, and pure endpoint squares remain. The four structural modifications are invertible, even when coefficient contractions are not.
\end{theorem}
\begin{proof}
The contextad is supplied by the first lemma and the independent occurrence/Fubini structure by the second. Its carrier operations are unchanged. The crossing at a content value of shape $s$ is $\widehat P_y=\boxtimes_{i\in E_s}P_{b_i}$ and its content grade is the original cartesian pullback of $Q$ along the context carrier projection. Content display, cartesian preservation, and Beck--Chevalley give the original body and grade maps.

The context-unit modification uses tensor units; content unit uses the singleton occurrence and cartesian content unit; context multiplication uses the new native Fubini lemma; content multiplication uses addressed leaf tensor reassociation and cartesian content accumulation. None uses an inverse coefficient contraction. The original ten law diagrams reduce to these native coefficient coherences and the unchanged cartesian uniqueness/mate-pasting comparisons. In their coefficient part, replace the former invertible-normalization argument by the directed indexed-source contractions proved above. Ancestor-dependent maps stay before their aggregations. Native arrows are not inverted, and carrier equality is not used to identify coefficient arrows.
\end{proof}

\section{A canonical conservative coefficient enlargement}
Assume the original coherent symmetric rig $\V$ is small in a fixed universe; a larger universe may be chosen explicitly. Let
\[
 \W=[\V,\mathbf{Set}]^{op},\qquad y:\V\longrightarrow\W,
 \quad y(V)=\V(V,-).
\]
Give the functor category the two Day convolutions induced by $\V^{op}$ and then take opposites. Garner--L\'opez Franco's directed rig convolution construction supplies the contraction comparisons. Both convolution operations remain symmetric monoidal.

The coherence has a direct element-level description before taking opposites. An expansion duplicates a presheaf element into the specified tagged branches, composes the original coefficient distributor, and takes the coend class. Either order of expanding a fixed finite expression duplicates each original element into the same branches. Cartesian diagonal coassociativity, counitality and symmetry, together with the original coefficient coherence and coend relations, equate the maps. Coend representatives generate the source; the argument applies also when a nullary expansion discards an element. Thus partition, mixed-variable, nullary, and tensor compatibility are available without invertibility. This is a use of ordinary set-valued presheaves; an arbitrary enriched replacement would need its own justification.

\begin{theorem}[Conservative native envelope]
Co-Yoneda induces a fully faithful functor on original augmented-family categories, strong for dependent accumulation and occurrence tensor. With the same content input it induces a strong locally fully faithful inclusion of the original mixed comma presentation and its interaction into the new one. It preserves the original actual pure endpoint squares. The coefficient category $\W$ is complete, but no completeness of $\Aug(\W)$ is asserted.
\end{theorem}
\begin{proof}
Co-Yoneda is fully faithful and strong for each Day tensor, including their units. On representables the directed comparisons are exactly the images of the old distributors. An old family is sent to the same finite carrier with coefficients $y(V_x)$ and augmentation obtained from the strong additive comparison and $y(\epsilon)$. For a fixed raw carrier map $f$, its aggregation hom-set becomes
\[
 \W\bigl(\bigoplus_{fx=x'}y(V_x),y(V'_{x'})\bigr)
 \cong\V\bigl(\bigoplus_{fx=x'}V_x,V'_{x'}\bigr).
\]
This bijection also reflects the augmentation equation, by faithfulness. Thus the whole native hom-set, including distinct maps with identical carriers and maps on empty fibres, is preserved and reflected.

The dependent and occurrence formulas consist of old tensor/sum structure, distributors, and actual input arrows. Strong preservation of these gives the stated coherent comparison isomorphisms. The content grade and its extension are unchanged. Hence in $H(A,B)=(e_AP\downarrow d_BQ)$ the context component is fully faithful, the content component identical, and the body equation unchanged. This proves local full faithfulness, not merely injectivity of output semantics. The same argument with context reindexing and unchanged content transport preserves endpoint squares and pure-middle balancing.

Completeness of $\W$ follows from cocompleteness of the covariant presheaf category. It does not imply limits in its augmented-family construction: native aggregation arrows cannot generally be decomposed into per-witness arrows.
\end{proof}

The enlargement can be strict. For the one-object terminal rig, $\W=\mathbf{Set}^{op}$, with both object-level convolutions given by Cartesian product of sets. An augmented coefficient family is then a finite family of pointed sets: the opposite augmentation chooses a point in their product. Native maps include contravariant pointed coefficient maps. A singleton outer family with a two-point coefficient is not isomorphic to any old family with singleton coefficients. This is a genuine new native grade, not a quotient or renaming. No physical interpretation of every such new grade is claimed.

\section{Scope and remaining work}
The retained-scope analysis explains why changing coefficients must be done carefully. Integrating one latent context across independent occurrence tensor requires a separate comonoid criterion; a new coefficient completion does not make all old random scopes copyable. Full faithfulness ensures that an absent original native map between old grades cannot appear merely through this enlargement.

The present result is a conservative interaction envelope, not the desired full feedback/concurrency completion. Shared joint apices require a further construction at the native augmented/context level. A native feedback operation must be supplied with its own admissibility, cut, and coefficient compatibility proof. Infinite effects need additional existence or convergence data. The old source is preserved, but arbitrary old limits are not asserted to remain universal against all new objects. No uniqueness, maximality, or free universal characterization of a feedback extension is inferred.

\begin{thebibliography}{9}
\bibitem{O1} \emph{General Independent Interaction Criterion}, project source, 1 November 2026. Original native depaug, depmap, tensormap and dependent Fubini; source blob \texttt{1315e4b4fca69e1899498ae831f4fe41b60b4d52}.
\url{https://github.com/eternal-november/eos-research-programmes/blob/main/projects/contextual-composition/interaction/independent-criterion/General_Independent_Interaction_Criterion.tex}.
\bibitem{GLF} R. Garner and I. L\'opez Franco, \emph{Commutativity}, Journal of Pure and Applied Algebra 220 (2016). Definition 42 and Proposition 43: directed rig comparisons and their convolution lift. \url{https://arxiv.org/pdf/1507.08710}.
\bibitem{Day} B. Day, \emph{On closed categories of functors}, Reports of the Midwest Category Seminar IV, Lecture Notes in Mathematics 137 (1970), 1--38. Convolution and strong Yoneda structure; the two-tensor directed extension used here is the preceding reference.
\end{thebibliography}
\end{document}
